% ============================================================================
% Chapter 3: Threat Model
% ============================================================================
\chapter{Threat Model}
\label{ch:threat_model}

This chapter defines the adversary model and security assumptions under which CBSR operates.

\section{Attacker Capabilities}

We assume a sophisticated attacker capable of:
\begin{itemize}
\item Executing arbitrary commands on compromised hosts---spawning processes, downloading
    payloads, and modifying file systems---without generating signature-based alerts.
\item Conducting multi-stage APT campaigns including initial access, privilege escalation,
    lateral movement, and data exfiltration.
\item Persisting in the target environment for extended periods (days to months).
\item Being \emph{aware of our monitoring presence} and deploying active evasion strategies.
\end{itemize}

\section{Fundamental Assumption}

We assume the attacker \textbf{cannot break the fundamental causal invariances of the operating
system without leaving traces in the provenance graph}. Any operation that creates a process,
reads or writes a file, or establishes a network connection \emph{must} be recorded in the
system audit log. This assumption holds for modern audit frameworks (Windows ETW, Linux auditd,
macOS Endpoint Security) that operate at the kernel level.

\section{Adaptive Adversary Strategies}

We consider three specific evasion strategies a sophisticated adversary might employ:

\textbf{Calibration Poisoning.}~A poisoner injects high-volume, anomalous but non-malicious-looking
traffic during the baseline training phase. This inflates the conformal calibration queue's score
distribution, widening detection threshold bounds and allowing subsequent stealthy attack steps to
bypass alerting. CBSR's RobustParCorr (MCD-based, tolerating up to 15\% contamination) insulates
the causal discovery phase; the Behavioral Pattern Reasoner's explicit validation-set training
further insulates the classification boundary from calibration inflation.

\textbf{Refit Blind Spots.}~When CBSR detects a concept drift change-point and refits its SCM,
a temporary adaptation window exists before calibration scores stabilize. An adversary can
intentionally trigger benign-looking drift (e.g., a software update) and immediately execute
malicious actions during this window. This vulnerability is bounded to approximately 1--2
conformal calibration cycles.

\textbf{SCM-Aware Mimicry.}~An attacker who learns the system SCM's structure can orchestrate
execution timings to simultaneously satisfy all $d \times \tau_{\max}$ causal constraints,
bypassing residual-based alerts. On DARPA OpTC ($d = 130$, $\tau_{\max} = 2$), this requires
matching 260 regression relationships while executing a multi-stage attack---a highly constrained
operational task. Moreover, real APT campaigns inject novel process--file edges that have no
causal parents in the baseline graph, triggering the Structural Consistency coordinate ($s_1$,
novelty count) regardless of residual behavior.

\section{Scope and Limitations}

CBSR focuses on defending against stealthy APTs where precise multi-constraint causal mimicry or
targeted poisoning is highly constrained. The framework is \emph{not} designed to defend against
attackers with kernel-level access who can tamper with the audit logging subsystem, or
side-channel attacks that leave no provenance traces. These limitations are discussed further in
Chapter~\ref{ch:conclusion}.
